\bookchapter{Carrier return and matter reception}
\label{sec:vii-carriers}

\subsection{How a closed object remains a source}

Finite closure does not isolate a material object from the substrate account. An admitted row continues to emit source support, and a composite must retain a rule for the support emitted by its connected configuration. The return route begins with the admitted object and ends with the source accepted by the accumulation field. It preserves the object as matter throughout the route.

The support decomposition established in Chapter~\ref{sec:vii-foundations} is
\begin{equation}
q_A\longmapsto
\left(T=\frac{q_A}{8},W=\frac{7q_A}{8}\right).
\end{equation}
The next operation is ledger compression, schematically
\begin{equation}
\mathcal C_{\rm ledger}(T,W,\text{source address})
\longmapsto r_s.
\end{equation}
Only after this operation does the spatial readout take the form
\begin{equation}
A(r)=\frac{r_s}{r}.
\end{equation}
The compression symbol identifies the interface supplied by the source construction. It is not an additional universal SI conversion derived in this chapter. The field readouts that consume $A(r)$ remain those of Volume I.\cite{viisama,viispine}

The order matters. Directly substituting $q_A$ for measured mass would omit the carrier and compression stages. Calling the released carrier another material constituent would instead count a channel of support as an added object. Both shortcuts change the source semantics even when they produce a familiar-looking scalar expression.

The present chapter follows this route through functional carrier and container roles, then examines the more detailed ATOM3D reception constructions. The early part concerns exact ledger distinctions and scoped catalog tests. The later part concerns finite information receipts, graph reconstruction and candidate support demand. The two parts are related, but neither supplies an automatic conversion from informational structure to physical length or energy.

\subsection{Conservation and the controls that identify its scope}

The source record QP092d reports conservation tests across seven particle-source rows, six macro-ledger rows and eighteen propagation rows. All 31 scoped rows pass their reconstruction checks, and three ledger scopes close. The reported maximum conservation relative error is approximately $2.98274132099\times10^{-19}$. The same record adds zero matter rows and rejects direct $q_A$-as-mass use on 24 of 24 tested rows.\cite{viireview}

The coexistence of successful conservation and failed direct mass identification is instructive. The exact decomposition concerns the sum of two support channels. It says nothing by itself about which downstream physical quantity should be identified with that sum. A construction can preserve its source quantity exactly while an attempted physical interpretation of the quantity fails.

The LC03 replay follows the longer matter-to-accumulation route and retains fifteen rejected wrong controls. The alternatives include missing carrier support, altered fractions, promotion of carriers to matter, vector/scalar substitution, mass-only source, omitted compression and direct use of support as measured mass. Each alternative changes a particular interface rather than merely perturbing an arbitrary output.

Some of these distinctions can be seen without a numerical test. Deleting $T$ while retaining $W$ leaves only $7q_A/8$, so the original support is no longer conserved. Other alternatives require the source's full typed structure. If both fractions were changed so that their sum remained one, the generic arithmetic conservation equation would still hold, but the specified one-eighth carrier relation to the structural domain would no longer be preserved. Conservation alone therefore does not select the source's fractions; the locked structural relation and its scoped controls are also part of the construction.

The tests should be read at this resolution. A rejected wrong fraction supports dependence on the specified domain relation. Its result has the scope of the specified carrier-transport mechanism. A successful replay establishes the mechanism it actually reconstructs, using the inputs and gates it records.

\subsection{Output fixed point and terminal closure}

Carrier and container roles are classified by source addresses, not by sorting their magnitudes. In CR281 the addresses have the form $h^a d^b$, with the primitive values associated with the binary and ternary structure. The classifier assigns terminal closure at address $162$; fixed-point behavior at $(a,b)=(2,1)$ or $a=0,b\ge3$; output behavior at $a=1,1\le b\le3$ or $(a,b)=(0,2)$; and OUT otherwise.\cite{viireview}

For the specific addresses used in the volume, this gives
\begin{align}
6,9,18,54&\longrightarrow\mathrm{OUTPUT},\\
12,27,81&\longrightarrow\mathrm{FIXED\ POINT},\\
162&\longrightarrow\mathrm{TERMINAL\ CLOSURE}.
\end{align}
The terminal case is evaluated as its declared special role. The list is a finite functional classification, not a monotone ordering by size. In particular, $12$ is fixed while $18$ is output. The source records six broken alternatives involving numerical ordering, parity, pure versus mixed powers, an unrestricted grade, swapped primitive addresses and inverted labels.

The role of $\Theta18$ illustrates why this classifier is needed. A hypothetical retained lift would produce
\begin{align}
18\left(1+\frac{18}{144}\right)
&=18\frac98\\
&=\frac{81}{4}.
\end{align}
The source does not seat that value as the carrier's material target row. The corresponding rest closure is $18-18=0$. It remains released output traffic. The same numeral can participate in a lift expression without thereby acquiring the right type for that expression's physical interpretation.

The archived QP101 audit provides an earlier inventory context: a 299-row V2 catalog with thirteen added roles, including five named carrier roles and eight lifted connector rows. That historical count is preserved as the scope of that audit. It does not replace the later 321-row catalog or imply that the thirteen roles are homogeneous. The current interpretation keeps X1, W8, W9, the lift-value alphabet, $\Theta18$ and F81 distinct.

\subsection{The non-row operator and the closed ledger}

The conjugation operator $\mathsf A$ transforms declared charged antecedents. In the compact laws its effect includes a sign- and address-dependent factor $C_A(s,p,d)$. The positive branch has a hard zero at $p=12,d=0$:
\begin{equation}
C_A(+)=\frac65\left(1-\frac{12}{12}\right)=0.
\end{equation}
The output is a consequence of the transform. The transform itself is not another particle-bearing row.

CR282 traces a historical A-field-carrier row that was retired as a footprint duplicate. Restoring that row would change the closed ledger from $162$ to $163$. The retained reading keeps the active operator concept while refusing the duplicate row count. This correction separates a useful operation from an inappropriate inventory representation of it.\cite{viireview}

A separate appeal source supports a route from contact through an X1 axis fee to W9, but its permanent direct test key is absent from the current registered source. The supported route is therefore identified as a separate source claim rather than silently assigned the key of the registered CR282 row-trace test. The direct relations from $A$ to contact and from $A$ to X1 remain open in that source account.

The custody distinction is scientifically important. Two documents can discuss related concepts while testing different propositions. A shared numerical identifier or a similar title is insufficient to transfer a result between them. The present manuscript retains the available route and the missing connection without creating a new causal edge between matter, activator and carrier.

\subsection{Cardinality closure without a geometric assignment}

From the completed relation $W9=8+1=9$, the source constructs the full-face value
\begin{equation}
F=9^2=81.
\end{equation}
Removing a non-row closure address leaves
\begin{equation}
F-1=80=(9-1)(9+1)=8\cdot10.
\end{equation}
The paired face ledger is $\mathcal L=2F=162$. These relations organize a useful cardinality contact: eighty active rows and one non-row address complete an eighty-one-position face.

Cardinality does not determine placement. CR283 tested two specified placements: a punctured $9\times9$ row map and a single-field ten-octet grouping. Both precommitted searches returned zero valid geometries. The result therefore has two parts. The arithmetic contact remains exact; the proposed row-to-face maps failed to supply the required geometry.\cite{viireview}

There is no need to discard the cardinality because the map failed, and no justification for claiming the map because the cardinality succeeded. A finite set can have the right number of members for a target without satisfying the target's adjacency, class or incidence constraints. Those constraints are additional mathematical content.

The N100-to-F81 audit makes the same distinction in a computation-facing setting. The reported source-native surface has 100 unique N100 keys, 72 rows inside the exact aggregate-incidence domain and an 81-site target with ten classes. It supplies zero source-authorized elementwise assignments; the contact lane is not executed. The resulting boundary prevents the missing map from being filled by fitting an outcome. It also keeps an executable grammar from rewriting the physical source tables it consumes.

\subsection{From ledgers to an ATOM3D scaffold}

The ATOM3D branch asks how typed matter information can be represented as exact incidence and, in selected cases, as an unscaled geometric readout. Its current authority distinguishes an informational construction from a calibrated physical nuclear model. The distinction is maintained even when a graph has a visually recognizable shape.\cite{viiatom}

The reported A1 O-16 scaffold installs eight identified protons, eight neutrons and eight electrons. It organizes 100 N100 roles into eight components, with 260 exact tensor relations, eight octahedral sites and 48 directed reciprocal actions. Its validation passes 14 of 14 checks. At that stage the scaffold assigns neither matter endpoints nor physical coordinates. The count of relations is a property of the installed ledger, not a measurement of nuclear geometry.

The hydrogen ladder then makes increasingly structured finite objects. H-1 is an empty-bridge control: one proton, no neutron, no inter-nucleon pair and no carrier event. The reported validation passes 22 of 22. H-2 installs the first reciprocal proton-neutron connector and calibrates it at $2.2245662$ MeV, with 29 of 29 validation checks. The energy here is a calibration. It is not an unanchored prediction of the preceding finite counts.

The H-3 and He-3 construction supplies reflected three-role paths and a frozen common prediction of $4.4491324$ MeV. It misses the two reference values by $4.03266380$ and $3.26890892$ MeV respectively. The source classification is falsification of that concept. The residual separates into a common part $3.65078636$ MeV and a signed half-difference of $0.38187744$ MeV. This decomposition organizes the recorded miss; it does not erase it.

A later signed-surface join matches the centered sign of the measured mirror relation while direct transfer of the surface magnitude fails. A target-visible discovery search then enumerates 13,641 expressions and finds a leading mirror half-split residual of $0.003918$ eV. The source assigns no confirmatory classification or installed physical law to that target-visible result. It is preserved as discovery, with its information access stated.

The sequence is useful because the successful structural pieces and the failed energy transfer can coexist. Reflection and exact incidence may be correctly implemented while a proposed calibration transfer is wrong. Treating a passing validator as proof of the physical prediction would lose that distinction: the validator can confirm that the frozen prediction was calculated correctly, including when the prediction disagrees with the measurement.

\subsection{He-4 and the difference between discovery and confirmation}

The reported He-4 pre-reveal construction forms an exact four-role N100 $K_4$. It contains six relations, six carrier cycles, twelve actions and six closures, with zero centered surface. Its self-mirror construction returns the ledger exactly after two reflections. These are structural results, and the source leaves them unclassified as physical-energy claims.

The frozen tetrahedral energy candidate combines six calibrated pair contributions and four common residual contributions:
\begin{align}
B_4={}&6(2.2245662)\\
&+4(3.65078636)\\
={}&27.95054264\ \mathrm{MeV}.
\end{align}
Against the recorded AME2020 value $28.2956624$ MeV, its miss is $0.34511976$ MeV. The exact-energy concept is classified as falsified. Within the frozen face-count comparison, four faces are nevertheless the unique best choice among zero through six, and that narrower result is recorded as strong contact.\cite{viiatom}

The target-visible T8C1 composition study then joins the eight-face Theta18 prism to the native partition alphabet. The inactive complement $\{2,3,4,6\}$ has squared fee $65/144$. Scaled by the A=3 mirror split, it yields a composition correction of $0.344750466666\ldots$ MeV. The resulting He-4 composition value is $28.295293106666\ldots$ MeV, leaving $369.293333\ldots$ eV. The source records no confirmatory classification because the target was visible.

The reduction in residual is substantial within the discovery calculation. Its epistemic status is still different from a frozen prediction tested against an unseen target. The chapter retains both facts. Neither dismissing the small residual nor presenting it as blind confirmation would preserve the actual record.

The same candidate extends to 24 balanced alpha multiples with $x=2,\ldots,25$, using
\begin{equation}
Z=N=2x,\qquad A_{\rm nuc}=4x,
\end{equation}
and a connector seam count $x-2$. The proposed readout is
\begin{equation}
B(x)=B_{4,\rm composition}\frac{5x-2}{4}.
\end{equation}
The reported relative RMS is $0.810343\%$ through Ca-40, $0.712682\%$ through Ni-56 and $3.757561\%$ for Zn-60 through Sn-100. The heavier drift remains in the record. A corrected validator passes 58 of 58 after an initial numeric-comparison defect left the first version at 54 of 58. These are candidate counts and readouts, not measured nuclear geometry or a general unbalanced binding law.

\subsection{Correcting the carbon and balanced-relation types}

The C12 type correction distinguishes three neutral seed cells from a whole-object coordinate. Three seeds give $6p+6n$. The aggregate quark coordinates are $u=d=18$, matching the scalar Theta value at the whole-object level. The correction does not establish a Theta multiplicity of three. The quantity $54=3\Theta$ remains an unassigned derived three-dimensional quantity in that record.

This is another instance of the central type discipline. A seed count and a whole-object coordinate can enter the same construction without being interchangeable. Multiplying a familiar scalar by the number of seeds does not establish that the object contains that many physical tensors. The corrected workbench retains the neutral seeds and places Theta18 at the whole-object level; the mistaken earlier artifacts remain preserved.

The balanced relation-emergence test reports closures $6,12,18,24,30$ before opening the comparison record and matches the balanced relation
\begin{equation}
I_{3D}=6N_{\rm seed}=\frac{3A_{\rm nuc}}2=3Z=u=d
\end{equation}
on five of five cases. Its classification is strong contact. The expression is scoped to the balanced seed construction; it is not a count law for arbitrary nuclear incidence.

The later loop-split construction preserves exact relation types and total $I_{3D}$ while redistributing closures. For C-12 it partitions eighteen as fourteen intra-seed and four inter-seed closures. The source assigns no new physical classification to this structural redistribution. Its importance is that connectivity can change while an aggregate count remains fixed. A count is therefore insufficient to identify the connected object.

Li-6 supplies a balanced non-alpha-multiple example. The target-barred N100 enumeration selects stable-matter depth-successor roles and emits nine relations: six $K_4$ addresses, one proton-neutron address and two same-role depth relations. All four connected balanced extensions return nine; 72 disconnected balanced controls return seven. Two same-charge candidates also return nine. Thus the numerical total alone is insufficient; the matter-role information is load-bearing. Exact incidence is installed, while action frames, physical geometry and binding readout for the added relations remain open.\cite{viiatom}

\subsection{The accumulation carrier and graph reconstruction}

AFC1 selects the separate A-field carrier for the local accumulation relay in sixteen of sixteen cases. This carrier is not the non-row conjugation operator $\mathsf A$. Maintaining the distinction prevents the accumulation field, an operator on charged antecedents and a transport role from being collapsed under the same letter.

AFC2 reconstructs a $3p+3n$ object from eighteen target-barred records. Its degree sequence is $4,4,3,3,2,2$, and its aggregate quark coordinates are $u=d=9$. Wrong-information controls can preserve the aggregate nine while losing the object. The result therefore tests more than a total: it relies on relation identities, endpoints and the information needed to reconstruct the graph.

AFC3 applies classical multidimensional scaling to the anonymous graph. The reported output has rank three, nonzero tetrahedral volume and four symmetries. The reveal identifies one full matter exchange, a proper 180-degree rotation about a common axis containing the three paired midpoints. Scrambled or disconnected ledgers fail the corresponding construction, while tampering and density-only alternatives change the shape.\cite{viiatom}

The coordinates are unscaled. Rank three establishes the dimension of the reported geometric representation under the selected construction, not a calibrated nuclear size. A proper rotation in that representation is a mathematical symmetry of the reconstructed object. It does not by itself identify a measured motion, a binding potential or a dynamical law. The source explicitly leaves physical length, $q_A$, binding and energy outside this result.

The distinction between exact incidence and displayed coordinates is useful. Incidence identifies which relations exist. A geometric embedding represents those relations using coordinates under a specified method. Physical calibration would be an additional map from that representation to measured units and observables. Preserving the three stages makes the finite reconstruction interpretable without overstating its physical reach.

\subsection{Relative support and complete delivery}

AFC4 joins the frozen relay and hologram ledger to the installed accumulation-conditioned support account. Matter-side $q_A$ remains source support, $q_A/8$ is the candidate carrier-support budget, and the A-field carrier remains non-matter. The reported gravity-only and velocity-only lanes agree when they have the same total accumulation.

With the inherited $A_0=1/(12\pi)$, the normalized demand is one baseline at $A=A_0$, $\pi$ baselines at $A=1/12$, $11\pi$ baselines at $A=11/12$, and $12\pi$ baselines at $A=1$. The source interprets the $11/12$ condition as the end of reliable old-object closure and the $A=1$ condition as a parent-road reset within that construction. These are inherited support hypotheses and relative outputs, not a newly derived SI-energy law.\cite{viiatom}

The finite delivery test enumerates all 512 reciprocal-delivery patterns of the nine-relation Li-6 ledger. Only complete delivery reproduces the exact AFC3 shape. Every single-relation loss changes it. Six of the nine single losses also destroy the full proton-neutron exchange; three preserve that exchange while deforming the object. Two missing relations can disconnect it. The reported independent validation passes 39 of 39.

This is a useful example of why both symmetry and object identity must be checked. Preserving one exchange symmetry does not guarantee preservation of the whole shape, and preserving a total count does not guarantee the required graph. Complete relation delivery supplies more information than those compressed diagnostics. The result is a finite, scoped statement about the reconstruction from its receipts.

It installs relative support demand, not joules and not a finite bounce-to-energy operator. The physical meaning of the extreme accumulation regime remains separately bounded. The chapter does not convert a successful information-delivery reconstruction into an experimental demonstration of matter transformation.

\subsection{Current partition families and local completion reports}

The current R4 partition study freezes 470 sum-12 families. Its reported minimum is 70 families over $\{3,9\}$ and $\{4,8\}$, while $\{2,4,6\}$ supplies a 100-family Q3-face witness with four closed routes. All 340 receipts distinguish 170 pairs, reverse and reproduce their sealed hashes; sixteen negative and two positive controls pass. The current authority retains strong contact for this finite informational result. Physical incidence, energy, Theta, training and the RH branch remain separate boundaries.\cite{viiatom}

The typed T63/W9 crosswalk supplies compatible account views:
\begin{align}
63+63+18&=144,\\
(63+9)+(63+9)&=144,\\
63+18&=81,\\
81+81&=162.
\end{align}
The source reports reconstruction across 48,620 custody partitions and 72 reciprocal receipts, together with 67,174,416 exchange cases. These are recorded finite checks. The type correction and Theta own-object fingerprint are retained, while a shared cross-layer identity receipt for the two-W9 intersection remains absent.

The local completion-report architecture expresses the roles in order: W8 and X1 complete W9; W9 supplies demand; local accumulation supplies conditions; $q_A$ supplies matter-side support; the A-field carrier transports it; and Theta reports completion at one local intersection. A distributed object is then described through its local reports. The architecture preserves these role assignments and their recorded physical transformations.

The scalar-parent relay study provides a further partial join. It identifies QP093A-0299 as the unique row that is both a closed scalar loop and a loop-closed-then-split object in the 321-row grammar. Its one-eighth and seven-eighth support fields agree with the generic support split. The source reports consistency across all 286 matter-allowed rows and failure of nine wrong controls, including a same-partition wrong parent. Exact shared receipts for the W9-to-scalar-parent and tensor-support-to-A-field-carrier edges remain open. Its classification is possibility, not an installed causal matter-generation law.

\subsection{The return that can be carried forward}

The finite return account is now explicit. An admitted material configuration emits typed support. Its carrier and retained channels conserve that support. Ledger compression supplies the interface to accumulation. Functional addresses distinguish output, fixed-point and terminal roles. Selected reception tests reconstruct exact graphs, unscaled geometric representations and relative support behavior from finite receipts.

The account preserves the failures that define its reach: no direct support-as-mass identity, no carrier promotion through a numerical lift, no duplicate operator row, no successful CR283 row-to-face geometry, no universal calibrated energy transfer, and no automatic physical interpretation of informational tensors. These limits do not remove the exact constructions. They identify the information that may be exported without alteration to the computation volume.
